\uuid{NgMU}
\exo7id{3767}
\titre{exo7 3767}
\auteur{quercia}
\organisation{exo7}
\datecreate{2010-03-11}
\isIndication{false}
\isCorrection{true}
\chapitre{Espace euclidien, espace normé}
\sousChapitre{Endomorphismes auto-adjoints}
\cpge{mpsi-mp2i.calcul-matriciel-systemes.e, pcsi.determinants.c, ptsi.determinants.a, mpsi-mp2i.matrices-applications-lineaires.c, pcsi.matrices-applications-lineaires.d, ptsi.matrices-applications-lineaires.d, bcpst1.matrices, bcpst1.geometrie.f, mp-mpi.reduction-endomorphismes.b, pc.reduction.a, psi.reduction-endomorphismes.a, pt.reduction-endomorphismes.a, mp-mpi.endomorphismes-euclidien.b, pc.endomorphismes-euclidiens.b, psi.endomorphismes-euclidien.b, pt.isometries-euclidien.b, bcpst2.valeurs-propres-vecteurs-propres.a, bcpst2.valeurs-propres-vecteurs-propres.b, bcpst2.produit-scalaire.c}
\module{Algèbre}
\niveau{L2}
\difficulte{}

\contenu{
\question{
Soient $a_1,\dots,a_n \in \R$ et $M = (a_ia_j) \in \mathcal{M}_n(\R)$.
Montrer que $M$ est diagonalisable et déterminer ses éléments propres.
}
\reponse{
Si tous les $a_i$ sont nuls, $M = 0$.\par
         Sinon, $M = C^tC  \Rightarrow  E_0 = C^\perp$ et $E_\nu = \text{vect}(C)$
         avec $\nu = \|C\|^2$.
}
}
